% !TeX root = ../thesis.tex
% ===================== ГЛАВА 1 =====================
\chapter{Математическая модель и численный метод}

\section{Уравнение Больцмана и модель твёрдых сфер}

Эволюция функции распределения молекул по скоростям $f(t,x,\xv)$ описывается кинетическим уравнением Больцмана
\begin{equation}
  \frac{\partial f}{\partial t} + \xi_x\frac{\partial f}{\partial x} = I(f,f),
  \label{eq:cbe}
\end{equation}
левая часть которого отвечает свободному переносу молекул, а правая~--- интегралу столкновений.
Для пространственно-однородного состояния (адвекции нет) в приближении молекулярного хаоса уравнение сводится к
уравнению релаксации
\begin{equation}
  \frac{\partial f}{\partial t} = I(f,f),\qquad
  I(f,f)=\int_{\mathbb{R}^3}\!\int_0^{2\pi}\!\int_0^{b_{\max}}\big(f'f'_1-f f_1\big)\,|\gv|\,b\dd b\dd\varepsilon\dd\xv_1,
  \label{eq:coll}
\end{equation}
где $\gv=\xv_1-\xv$~--- относительная скорость сталкивающейся пары, $b$~--- прицельное расстояние, $\varepsilon$~---
азимутальный угол, а штрихом обозначены скорости после столкновения.
Используется модель твёрдых сфер диаметра $d$, для которой $b\in[0,d]$ ($b_{\max}=d$), а угол отклонения относительной
скорости равен $\theta=2\arccos(b/b_{\max})$.
Равновесным решением~\eqref{eq:coll} является распределение Максвелла
\begin{equation}
  f_M(\xv)=n\left(\frac{1}{2\pi T}\right)^{3/2}\exp\!\left(-\frac{(\xv-\mathbf u)^2}{2T}\right),
  \label{eq:Maxwell}
\end{equation}
где $\mathbf u$~--- средняя (гидродинамическая) скорость газа, $n$~--- концентрация, $T$~--- температура.
Это распределение обращает интеграл столкновений в ноль.

Основная трудность численного решения~--- именно интеграл столкновений: он нелинеен по $f$ и имеет высокую кратность
(восемь измерений в полном пространстве параметров пары), тогда как левая часть уравнения~--- линейный оператор
переноса.
Это определяет общую стратегию работы: оператор столкновений вычисляется консервативным проекционным методом
(раздел~1.4), а в неоднородной задаче он отделяется от оператора переноса схемой расщепления (глава~3).
Интеграл столкновений обладает важными свойствами: он сохраняет суммарные число частиц, импульс и энергию (аддитивные
инварианты столкновения), а также обеспечивает невозрастание $H$-функции $H = \int f(\xv) \ln f(\xv) \dd\xv$
($H$-теорема)~\cite{LL10,Kogan}, причём $H$ достигает минимума именно на распределении Максвелла.
Сохранение этих свойств на уровне разностной схемы и является главным требованием к численному методу.

\section{Обезразмеривание}

Перед записью разностной схемы перейдём к безразмерным переменным.
За масштабы примем температуру опорного газа $T_0$ (набегающего потока в задаче о поршне; холодной компоненты в задаче о
смеси) и концентрацию $n_0$, диаметр молекул $d$, тепловую скорость $\xi_0=\sqrt{RT_0}=\sqrt{k_B T_0/m}$ (здесь
$R=k_B/m$~--- удельная газовая постоянная) и среднюю длину свободного пробега
\begin{equation}
  \lambda=\frac{1}{\sqrt{2}\,\pi d^2 n_0}.
\end{equation}
Безразмерные скорость, время и функция распределения вводятся как $\xv\to\xv/\xi_0$, $t\to t/\tau_\lambda$, $f\to
f\,\xi_0^3/n_0$, где \emph{характерное время между столкновениями} $\tau_\lambda=\lambda/\xi_0$ принято за единицу
времени; за несколько $\tau_\lambda$ функция распределения приходит к равновесию.
В безразмерных переменных перед интегралом появляется множитель $1/Kn$, имеющий смысл обратного числа Кнудсена.
Шаг по времени схемы (приращение безразмерного времени $t$, измеряемого в единицах $\tau_\lambda$) всюду $\tau=\Delta
t=0{,}02$; само $t\to t/\tau_\lambda$~--- это время в единицах $\tau_\lambda$, а $\tau$~--- его малый шаг.

\section{Дискретизация пространства скоростей}

В численной реализации вводится равномерная сетка скоростей
\begin{equation}
  \xi_{\alpha_i}=-\xi_{\mathrm{cut}}+(\alpha_i-0{,}5)\,\Delta\xi_i,\qquad
  \Delta\xi_i=\frac{2\xi_{\mathrm{cut}}}{N_i},\qquad \alpha_i=1,\dots,N_i,
\end{equation}
(аналогично по осям $y$ и $z$).
Пространство скоростей разбивается на одинаковые ячейки-кубики объёмом $\Delta V=\Delta\xi_x\Delta\xi_y\Delta\xi_z$;
функция распределения считается постоянной в пределах ячейки и относится к её центру.
Область ограничена сферой обрезания $|\xv|\le\xi_{\mathrm{cut}}$: функция распределения быстро спадает с ростом
скорости, и вкладом скоростей в несколько среднеквадратичных можно пренебречь.
На дискретной сетке макропараметры вычисляются конечными суммами
\begin{equation}
  n=\sum_\alpha f_\alpha\,\Delta V,\quad
  \mathbf u=\frac1n\sum_\alpha \xv_\alpha f_\alpha\,\Delta V,\quad
  T=\frac{1}{3n}\sum_\alpha (\xv_\alpha-\mathbf u)^2 f_\alpha\,\Delta V,
\end{equation}
\begin{equation}
  T_{xx}=\frac1n\sum_\alpha (\xi_{x,\alpha}-u_x)^2 f_\alpha\,\Delta V,\qquad
  H=\sum_\alpha f_\alpha\ln f_\alpha\,\Delta V,\quad S=-H.
\end{equation}

\section{Консервативный проекционный метод Черемисина}

Производящий интеграл~\eqref{eq:coll} является \emph{восьмикратным} в полном пространстве параметров столкновения: три
компоненты скорости узла $\xv_\alpha$, три компоненты скорости партнёра $\xv_\beta$, прицельное расстояние $b$ и
азимутальный угол $\varepsilon$.
Для каждой пары скоростей вычисляются пост-столкновительные скорости: относительная скорость $\gv$ поворачивается на
угол $\theta$ в плоскости $\varepsilon$, давая $\gv'$, после чего
\begin{equation}
  \xv'_\alpha=\tfrac12(\xv_\alpha+\xv_\beta)-\tfrac12\gv',\qquad
  \xv'_\beta=\tfrac12(\xv_\alpha+\xv_\beta)+\tfrac12\gv'.
\end{equation}
Как правило, $\xv'_\alpha$ не попадает в узел сетки, поэтому результат \emph{проецируется} на пару узлов $\lambda$ и
$\lambda+s$ с весами $1-r_\nu$ и $r_\nu$; симметрично $\xv'_\beta$ проецируется на пару $\mu$ и $\mu-s$.
Условие
\begin{equation}
  \mu=\alpha+\beta-\lambda
\end{equation}
обеспечивает точное сохранение импульса, а коэффициент $r_\nu\in[0,1]$ выбирается из условия сохранения энергии (для
обеих проекционных пар совместно)
\begin{equation}
  (1-r_{\nu})\big(E_{\lambda}+E_{\mu}\big)+r_\nu\big(E_{\lambda+s}+E_{\mu-s}\big)=|\xv_\alpha|^2+|\xv_\beta|^2,
\end{equation}
где $E_k=|\xv_k|^2$~--- кинетическая энергия в узле $k$.
Это и есть проецирование непрерывного отклика интеграла столкновений на дискретные узлы сетки~--- ключевая идея метода.
Дискретный вклад в узлы записывается через \emph{геометрическую} (логарифмическую) интерполяцию
\begin{equation}
  \Omega_\nu=\Big[\big(f_\lambda f_\mu\big)^{1-r_\nu}\big(f_{\lambda+s}f_{\mu-s}\big)^{r_\nu}
              -f_\alpha f_\beta\Big]\,|\gv_\nu|,\qquad \Delta_\nu=C\,\Omega_\nu,
  \label{eq:omega}
\end{equation}
с константой $C=\dfrac{b_{\max}^2\,V_{\mathrm{sph}}\,N_0}{4\sqrt{2}\,N_\nu}\,\tau$, где множитель $\tau$~--- шаг по
времени (схема явная по Эйлеру $f^{n+1}=f^n+\tau\,I$, поэтому отклик интеграла столкновений умножается на $\tau$),
$V_{\mathrm{sph}}=\tfrac43\pi\xi_{\mathrm{cut}}^3$~--- объём сферы обрезания, $N_0$~--- число узлов сетки внутри неё,
$N_\nu$~--- число точек сетки Коробова, попавших одновременно в сферу по обеим скоростям.

Полное обновление функции распределения за один шаг по времени имеет явный вид
\begin{equation}
  f_\gamma^{n+1}=f_\gamma^{n}+\sum_{\nu}\Delta_\nu\Big(\delta_{\gamma\alpha_\nu}+\delta_{\gamma\beta_\nu}
  -(1-r_\nu)\big(\delta_{\gamma\lambda_\nu}+\delta_{\gamma\mu_\nu}\big)
  -r_\nu\big(\delta_{\gamma,\lambda_\nu+s_\nu}+\delta_{\gamma,\mu_\nu-s_\nu}\big)\Big),
\end{equation}
где $\delta$~--- символ Кронекера, распределяющий вклад $\Delta_\nu$ каждого столкновения по шести участвующим узлам
($\alpha,\beta,\lambda,\mu$ и двум смещённым).
Применяется схема непрерывного счёта: обновления вносятся внутри цикла по $\nu$, а не накапливаются после полного
суммирования, что повышает устойчивость при увеличенном шаге $\tau$.
Геометрическая форма~\eqref{eq:omega} (а не линейная) гарантирует положительность $f$ и невозрастание $H$-функции
($H$-теорема).

\section{Сетки Коробова}

Для оценки восьмикратного интеграла каждая точка цикла $\nu$ задаёт сразу все восемь параметров столкновения из одной
8-мерной теоретико-числовой сетки Коробова~\cite{Korobov} на простом числе $p$:
\begin{equation}
  u_j^{(\nu)}=\left\{\frac{a_j\,\nu}{p}+\delta_j\right\},\qquad j=1,\dots,8,
\end{equation}
где $\{\cdot\}$~--- дробная часть, $a_j$~--- оптимальные коэффициенты, а случайный сдвиг $\delta_j$ обновляется на
каждом шаге по времени.
Из координат восстанавливаются физические параметры: $\varepsilon=2\pi u_1$, $b^2=b_{\max}^2 u_2$; компоненты первой
скорости $\xv_\alpha$ берутся из $u_3,u_4,u_5$, второй $\xv_\beta$~--- из $u_6,u_7,u_8$ линейным отображением
$[0,1]\to[-\xi_{\mathrm{cut}},\xi_{\mathrm{cut}}]$, то есть $\xi=(2u-1)\,\xi_{\mathrm{cut}}$, что равномерно покрывает
куб скоростей.
Точки, у которых хотя бы одна скорость выходит за сферу обрезания, исключаются из суммы методом отбора.
Такая кубатура даёт погрешность $O\!\left((\ln p)^8/p\right)$, существенно меньшую, чем $O\left(1/\sqrt{N}\right)$
метода Монте-Карло при том же числе узлов.
В расчётах использованы простые числа $p$ из таблиц Пособия~\cite{Tcher} ($p=50021$ и крупнее).

\section{Контроль положительности и законы сохранения}

Обновления применяются внутри цикла по точкам Коробова.
Если очередное обновление сделало бы какое-либо из значений $f$ в участвующих узлах отрицательным, текущая точка сетки
Коробова не даёт вклада в сумму.
Для «спокойных» однородных задач релаксации этот порог жёсткий ($\sim10^{-4}$, и на практике доля пропусков нулевая); на
существенно неравновесном фронте ударной волны (глава~3) допустима доля порядка процента, поскольку пропуск точки
Коробова не нарушает законов сохранения, а лишь не вносит вклад в сумму.
В отличие от методов дискретных скоростей (Бродвелл, Кабаннес), где разрешены только столкновения, точно попадающие в
узлы сетки, и потому сужается сам набор столкновений, здесь пропуск~--- это лишь статистическая потеря одной точки
кубатуры из многих: набор разрешённых столкновений и инварианты не меняются, физика не искажается.
Превышение порога означает необходимость увеличить $p$ или уменьшить $\tau$.
Кроме того, при чётном по скоростям начальном условии оператор столкновений сохраняет эту симметрию, что позволяет
хранить функцию распределения в виде четверти (либо восьмой части) массива и экономить память.
Проекционная конструкция обеспечивает точное (машинное) сохранение массы, импульса и энергии, что используется далее как
критерий корректности расчётов.
